% compile with XeLaTeX
\documentclass[dvipsnames,mathserif]{beamer}
\setbeamertemplate{footline}[frame number]
\setbeamercolor{footline}{fg=black}
\setbeamerfont{footline}{series=\bfseries}
\usepackage{tikz}

%\usetheme{Frankfurt}%1
\usetheme{Darmstadt}%1

% for RTL liste
\makeatletter
\newcommand{\RTListe}{\raggedleft\rightskip\leftm}
\newcommand{\leftm}{\@totalleftmargin}
\makeatother

% RTL frame title
\setbeamertemplate{frametitle}
{\vspace*{-1mm}
  \nointerlineskip
    \begin{beamercolorbox}[sep=0.3cm,ht=2.2em,wd=\paperwidth]{frametitle}
        \vbox{}\vskip-2ex%
        \strut\hskip1ex\insertframetitle\strut
        \vskip-0.8ex%
    \end{beamercolorbox}
}
% align subsection in toc
\makeatletter
\setbeamertemplate{subsection in toc}
{\leavevmode\rightskip=5ex%
  \llap{\raise0.1ex\beamer@usesphere{subsection number projected}{bigsphere}\kern1ex}%
  \inserttocsubsection\par%
}
\makeatother

% RTL triangle for itemize
\setbeamertemplate{itemize item}{\scriptsize\raise1.25pt\hbox{\donotcoloroutermaths$\blacktriangleleft$}} 

%\setbeamertemplate{itemize item}{\rule{4pt}{4pt}}

\defbeamertemplate{enumerate item}{square2}
{\LR{
    %
    \hbox{%
    \usebeamerfont*{item projected}%
    \usebeamercolor[bg]{item projected}%
    \vrule width2.25ex height1.85ex depth.4ex%
    \hskip-2.25ex%
    \hbox to2.25ex{%
      \hfil%
      {\color{fg}\insertenumlabel}%
      \hfil}%
  }%
}}

\setbeamertemplate{enumerate item}[square2]
\setbeamertemplate{navigation symbols}{}


\titlegraphic { 
\begin{tikzpicture}[overlay,remember picture, opacity=0.1,]
\node[] at (0, 2.9){
    \includegraphics[width=0.63\textwidth]{Logo_uzad.png}
};\end{tikzpicture}}
\setbeamertemplate{caption}[numbered]
\begin{document}

\rightskip\rightmargin
\title{Control and stabilization of delayed systems  }
\author{ \Large \textbf{Xxxx YYYYYY} }
\institute{\large\textbf{Master Automatic control and systems }\\
---------\\
Faculty Sciences and Technology\\
---------\\
University of Djelfa}
\footnotesize{\date{\today }


\begin{frame}
\maketitle
\end{frame}

\begin{frame}{ Summery}
\footnotesize \tableofcontents
\end{frame}

\section{Introduction}
\begin{frame}
\large Motivations:
\begin{itemize}
    \item Most systems are nonlinear
    
    \item Delay complicates the system analysis
   
    \item Delay can lead to the system instability
    
    \item ...    
\end{itemize}
\end{frame}
\section{Modeling}
\begin{frame}
\footnotesize
The TS fuzzy model can be justified by:
\begin{itemize}
    \item Its simplicity
    \item Uncertainties  
    \item Its acceptable accuracy
    \item ....
\end{itemize}
  \pause
\quad \\
The system can be modeled as: 
\begin{equation*}
\left\{\begin{array}{l}
{ }^C D^\alpha x(t)=f(x(t),x(t-\tau(t)),u(t)),\, t \geq 0, \\
x(s)=\varphi(s),\, s \in[-\tau, 0]
\end{array}\right.
\end{equation*}
$x(t) \in \Re ^{n}$ the system state\\
$u(t) \in \Re ^{m}$ the control vector \\
$\tau$: the delay
\end{frame}

\section{Lemmas}
    \begin{frame}
    \footnotesize
       \begin{block}{Lemma 1}
        \begin{equation*}
           I(x,t)=\int_{a(t)}^{b(t)} f(x,t) \,dx \  \ ,  \ a(t) \And b(t)<\infty
         \end{equation*}
              with 
              $a(t)$ و $b(t)$  
and               $f(x,t)$
                
         \begin{equation*}
   \frac{d(I(x,t))}{dt}=\frac{db(t)}{dt}f(b(t),t)-\frac{da(t)}{dt}f(a(t),t)+\int_{a(t)}^{b(t)} \frac{\partial f(x,t)}{\partial t} \,dx  
   \end{equation*}
       \end{block}
\begin{block}{Lemma 2}
        \begin{equation*}
           I(x,t)=\int_{a(t)}^{b(t)} f(x,t) \,dx \  \ ,  \ a(t) \And b(t)<\infty
         \end{equation*}
              with 
              $a(t)$ و $b(t)$  
and               $f(x,t)$
                
         \begin{equation*}
   \frac{d(I(x,t))}{dt}=\frac{db(t)}{dt}f(b(t),t)-\frac{da(t)}{dt}f(a(t),t)+\int_{a(t)}^{b(t)} \frac{\partial f(x,t)}{\partial t} \,dx  
   \end{equation*}
       \end{block}
\end{frame}    %.............................................................


\footnotesize

\section{Control without delay}
\begin{frame}
\footnotesize
\begin{block}{Theorem}
Let's consider:
\begin{equation*}
\left\{\begin{array}{l}
{\Omega_{ii}}< 0,\   (i=1,2,...r) , \\
{\Omega_{ij}}+{\Omega_{ji}}<0,\ i<j ,\ i,j=1,2,...r
\end{array}\right.
\end{equation*}
\end{block}
with:
\tiny
\begin{equation*}
\Omega_{ij}= \left[\begin{array}{cccc}
PA_i+A_i^TP^T+PB_iK_j+K_j^TB_i^TP^T+Q& PA_{di}& A_i^TN^T+K_j^TB_i^TN^T & 0 \\ *& -(1-\mu)Q& A_{di}^TN^T & 0\\ *&*&\tau^2R-N-N^T & 0\\ * &* &* &-R 
   \end{array}\right]
\end{equation*}
\end{frame}
\section{Proposed controller}
\footnotesize
\begin{frame}
We consider:
$u(t)=\sum_{i=1}^{r} h_{i}(\theta(t))[K_{i}x(t)+K_{di}x(t-\tau(t)]$ 
\begin{block}{Theorem }

\begin{equation*}
\left\{\begin{array}{l}
  \Bar{\Omega_{ii}}< 0,\ i=1,2,...r \\
    \Bar{\Omega_{ij}}+\Bar{\Omega_{ji}}<0,\ i<j ,\ i,j=1,2,...r
\end{array}\right.
\end{equation*}
\end{block}
where:\\
$ \Bar{\Omega_{ij}}=\left[\begin{array}{cccc}
A_iX+X^TA_i^T+B_iY_j+Y_j^TB_i^T+\Bar{Q}& A_{di}X+B_iY_{dj}& \epsilon X^T A_i^T+\epsilon Y_j^TB_i^T & 0 \\ *& _(1-\mu)\Bar{Q}& \epsilon X^T A_{di}^T+\epsilon Y_{di}^TB_i^T & 0\\ *&*&\tau^2\Bar{R}-\epsilon X-\epsilon X^T & 0\\ *&*&*&-\Bar{R}
   \end{array}\right]$
 $K_j=Y_jX^{-1}$ و $K_{dj}=Y_{dj}X^{-1}(j=1,2,...r)$
\end{frame}  

\section{Simulation results}

\begin{frame}
\footnotesize
System behavior without controller
\begin{figure}[H]
  \centering
  \includegraphics[width=0.75\textwidth]{x_1.eps}
  \caption{System state}
  \end{figure}
System unstable
\end{frame}

\begin{frame}
System behavior with controller

\begin{figure}[H]
    \centering
    \includegraphics[width=0.75\textwidth]{kuc.eps}
    \caption{Control signal } 
    \label{fig3.3}
\end{figure}
\end{frame}
\begin{frame}
System behavior with controller
\begin{figure}[H]
    \centering
    \includegraphics[width=0.75\textwidth]{kstates.eps}
    \caption{System state} 
    \label{fig3.4}
\end{figure} 

The system is stable but it needs more enhancement    
\end{frame}


\begin{frame}       
System behavior with the proposed controller
\begin{figure}[H]
    \centering
    \includegraphics[width=0.75\textwidth]{2uc.eps}
    \caption{Proposed controller signal } 
    \label{fig3.5}
\end{figure}
\end{frame}
\begin{frame}
System behavior with the proposed controller

\begin{figure}[H]
    \centering
    \includegraphics[width=0.75\textwidth]{2states.eps}
    \caption{System state with the proposed controller } 
    \label{fig3.6}
\end{figure} 

\textcolor{blue}{ Stability + better performance}
 
\end{frame}

\begin{frame}
\footnotesize
Quantification of the comparative study
\begin{table}[h]
    \centering
    \begin{tabular}{|*{4}{c|}}
\hline 
  & Classical controller & Proposed controller & Enhancement rate \\
\hline 
Settling time & $26$  &  $20$ & 23 \% \\
\hline 
Pic to pic $x_3$  &$1.36$ & $1.16 $ &15 \% \\
\hline
$\int\limits_0^{ts}(x_3^2)dt $
 & $2.5540$ & $0.7771$& 70 \% \\
\hline
$ \int\limits_0^{ts}(u^2)dt$ & $12.8476$ & $7.1868$& 40 \% \\
\hline
\end{tabular}
    \caption{Quantification of the comparative study}
\end{table}
We remark that
\begin{itemize}
    \item \textcolor{blue}{\textbf{Enhancement of the settling time of $23\%$}}
    \item \textcolor{blue}{\textbf{ Reduction of the control energy by $40\%$}}
    \item \textcolor{blue}{\textbf{ Overall enhancement by  $70\%$ }}
    \item \textcolor{blue}{\textbf{ Pic to pic reduction by  $15\%$}}
\end{itemize}
\end{frame}

\section{Conclusion and Perspectives}
\begin{frame}
\footnotesize
We conclude that:
\begin{block}{Conclusion}
\begin{itemize}
    \item  \textcolor{blue}{\textbf{Lyapunov method efficiency.}}
        \item \textcolor{blue}{\textbf{Proposed controller leads to better performance.}}
        \item \textcolor{blue}{\textbf{Delayed controller enhances the performance.}}
    \item \textcolor{blue}{\textbf{Proposed approach allows reduction of the control energy.}}
\end{itemize}
\end{block}
As perspectives we propose:
\begin{block}{perspectives}
\begin{itemize}
    \item Perspective 1.
    \item Perspective 2.
    \item Perspective 3.
\end{itemize}
\end{block}
\end{frame}


\end{document}









